\newcommand{\conone}{
In this assignment, a Python program was developed to read marks from three different subjects and calculate their average. The program demonstrates the practical use of basic Python concepts such as variables, input and output functions, and arithmetic operators.

By taking user input and applying the average formula:

\[
\text{Average} = \frac{\text{Subject1} + \text{Subject2} + \text{Subject3}}{3}
\]

the program shows how Python can be used to solve simple real-life problems efficiently and accurately.

This exercise improves logical thinking and helps in understanding how data is processed within a program. It also strengthens the foundation of programming concepts required for building more advanced applications in the future.

Overall, the objective of calculating the average marks was successfully achieved, highlighting the simplicity and effectiveness of Python programming.

}

\newcommand{\contwo}{
In this assignment, a Python script was developed to convert weight from kilograms to pounds. The program takes the weight in kilograms as input from the user and applies the standard conversion formula to calculate the equivalent weight in pounds.

\[
\text{Pounds} = \text{Kilograms} \times 2.20462
\]

The implementation demonstrates the use of Python variables, user input, arithmetic operations, and output statements. It provides a simple and efficient way to perform unit conversion accurately.

Through this exercise, a better understanding of basic programming concepts and mathematical computations in Python was achieved. The program successfully met its objective by converting weight values correctly and displaying the results in a user-friendly format.

Overall, this assignment highlights the effectiveness of Python for solving practical problems involving calculations and unit conversions.
}

\newcommand{\conthree}{
In this assignment, a Python program was developed to calculate the surface area of a rectangular prism using the lengths of its three sides as input. The program accepts three integer values representing the dimensions of the prism and applies the surface area formula:

\[
\text{Surface Area} = 2ab + 2bc + 2ca
\]

where \(a\), \(b\), and \(c\) are the lengths of the three sides.

The implementation demonstrates the use of user input, variables, arithmetic operations, and output functions in Python. By performing the necessary calculations automatically, the program provides an accurate and efficient method for determining the surface area of a prism.

This exercise helped in strengthening the understanding of mathematical computations and problem-solving techniques using \\Python. It also provided practical experience in translating mathematical formulas into executable code.
}

\newcommand{\confour}{
In this assignment, a Python program was designed to calculate the speed of a plane using a given distance and travel time. The program uses the standard speed formula:

\[
\text{Speed} = \frac{\text{Distance}}{\text{Time}}
\]

where the distance is 395,000 meters and the time taken is 9,000 seconds.

By substituting the given values, the program calculates the speed as:

\[
\text{Speed} = \frac{395000}{9000} \approx 43.89 \text{ m/s}
\]

The implementation demonstrates the use of variables, arithmetic operations, and output statements in Python. It provides an accurate and efficient way to perform calculations related to motion and speed.

This exercise helped in understanding how mathematical formulas can be translated into Python code and applied to solve real-world problems. It also strengthened basic programming skills and logical thinking.
}

\newcommand{\confive}{
In this assignment, a Python program was developed to calculate the time required to empty a swimming pool using a pump with a known flow rate. The pool dimensions are given as 12 m × 7 m × 2 m.

First, the volume of the swimming pool is calculated using the formula:

\[
\text{Volume} = \text{Length} \times \text{Width} \times \text{Height}
\]

\[
\text{Volume} = 12 \times 7 \times 2 = 168 \text{ cubic meters}
\]

The pump flow rate is given as 17 cubic meters per hour. The time required to empty the pool is calculated using:

\[
\text{Time} = \frac{\text{Volume}}{\text{Flow Rate}} = \frac{168}{17} \approx 9.88 \text{ hours}
\]
}

\newcommand{\consix}{
In this assignment, a Python script was developed to convert temperature from Celsius to Fahrenheit. The program takes a floating-point value as input from the user and applies the standard conversion formula:

\[
\text{Fahrenheit} = (\text{Celsius} \times \frac{9}{5}) + 32
\]

The implementation demonstrates the use of variables, user input, type conversion, and arithmetic operations in Python.
}

\newcommand{\conseven}{
In this assignment, a Python program was developed to compute the total number of seconds in one full day. The program uses basic time conversion concepts to perform the calculation.

We know that:
\[
1 \text{ day} = 24 \text{ hours}, \quad 1 \text{ hour} = 60 \text{ minutes}, \quad 1 \text{ minute} = 60 \text{ seconds}
\]

So, the total number of seconds in a day is calculated as:
\[
\text{Seconds in a day} = 24 \times 60 \times 60 = 86400
\]

The program demonstrates the use of arithmetic operations and constant values in Python. It provides a simple and efficient way to perform time unit conversions without manual calculation.

This exercise helped in understanding how basic mathematical relationships between time units can be implemented in programming. It also improved logical thinking and coding clarity.

Overall, the objective of computing the total number of seconds in a full day was successfully achieved using Python.
}

\newcommand{\coneight}{
In this assignment, a Python program was developed to calculate the acceleration of a car that starts from rest and reaches a final velocity in a given time. The initial velocity is 0 m/s, the final velocity is 10 m/s, and the time taken is 20 seconds.

The acceleration is calculated using the standard formula:

\[
\text{Acceleration} = \frac{v_{\text{final}} - v_{\text{initial}}}{\text{time}}
\]

Substituting the given values:

\[
\text{Acceleration} = \frac{10 - 0}{20} = 0.5 \text{ m/s}^2
\]

The program demonstrates the use of variables, arithmetic operations, and basic physics concepts implemented in Python. It provides an efficient way to solve motion-related problems using simple computational logic.

This exercise helped in understanding how real-world physics formulas can be translated into Python code for automated calculation. It also improved problem-solving and programming skills.
}

\newcommand{\connine}{
In this assignment, a Python program was developed to generate and display Pascal’s Triangle up to a given number of rows. The program uses nested loops and basic arithmetic logic to construct each row based on the sum of adjacent elements from the previous row.

The implementation demonstrates the effective use of loops, list operations, and pattern generation techniques in Python. It helps in understanding how mathematical patterns can be translated into programmatic logic.

Through this exercise, the concept of combinatorial structures and sequence generation was better understood. It also improved logical thinking and coding skills by applying iterative methods to build structured output.

Overall, the objective of displaying Pascal’s Triangle was successfully achieved using Python programming, highlighting its usefulness in mathematical problem-solving and pattern generation.
}

\newcommand{\conten}{
In this assignment, a Python program was developed to generate and display a number pattern using nested loops. The pattern was constructed such that each row follows a specific rule based on the row number, combining repeated digits and descending sequences.

The implementation demonstrates the use of nested loops, conditional logic, and output formatting in Python. It highlights how repetitive structures and patterns can be efficiently generated using programming constructs.

Through this exercise, a better understanding of loop control and pattern generation techniques was achieved. It also improved logical thinking and problem-solving skills by translating a structured pattern into code.

Overall, the objective of displaying the required number pattern was successfully achieved using Python programming, showcasing the effectiveness of nested loops in pattern-based problems.
}

\newcommand{\coneleven}{
In this assignment, a Python program was developed to solve the given problem efficiently using basic programming concepts. The program takes the required inputs and performs the necessary calculations using arithmetic operations.

The logic of the program is based on simple mathematical reasoning, which is implemented using Python syntax in a clear and structured manner. It helps in understanding how real-world problems can be translated into computational steps.
}

\newcommand{\contwelve}{
In this assignment, a Python program was developed to compute the sum of two different mathematical series using iterative logic. The program takes the value of \(n\) as input and calculates the result step by step using loop structures.

For the first series:
\[
1 + \frac{1}{2} + \frac{1}{3} + \cdots + \frac{1}{n}
\]
the program uses a loop to add the reciprocal of each integer from 1 to \(n\), generating the harmonic series sum.

For the second series:
\[
\frac{1}{1} + \frac{2^2}{2} + \frac{3^3}{3} + \cdots + \frac{n^n}{n}
\]
the program calculates each term using exponentiation and division, and accumulates the total sum using iteration.

The implementation demonstrates the use of loops, arithmetic operations, and mathematical functions in Python. It also helps in understanding how series-based mathematical problems can be converted into computational logic.
}
\newcommand{\conthirteen}{
In this assignment, a Python program was developed to calculate the depreciation value of an asset based on its purchase value, service years, and depreciation rate. The program takes input values such as the initial amount (purchase value), number of years of service, and the depreciation rate.

The depreciation is calculated using a standard mathematical approach, where the value of the asset decreases gradually over time according to the given rate. The program applies arithmetic operations and iterative logic (if required) to compute the final depreciated value.

The implementation demonstrates the use of variables, user input, and mathematical calculations in Python. It helps in understanding how financial concepts like depreciation can be modeled using programming.

This exercise improved logical thinking and provided practical knowledge of applying programming to real-world financial problems.


}
\newcommand{\confourteen}{
In this assignment, a Python program was developed to accept a string input from the user and display the same string after removing all vowels. The program processes each character of the string and filters out vowels such as a, e, i, o, and u (in both uppercase and lowercase).

The implementation demonstrates the use of string manipulation, loops, and conditional statements in Python. It shows how each character of a string can be checked and selectively included or excluded based on a given condition.

Through this exercise, a better understanding of string handling and character filtering techniques was achieved. It also improved logical thinking and programming skills by applying conditions to process textual data.

Overall, the objective of removing vowels from a given string was successfully achieved using Python programming, highlighting the effectiveness of Python in string processing tasks.
}
\newcommand{\confifteen}{
In this assignment, a Python function was developed to insert one string into the middle of another string. The program takes two strings as input: the main string and the string to be inserted.

The function calculates the middle position of the main string using string length and then splits the string into two halves. The new string is formed by concatenating the first half, the inserted string, and the second half.

The implementation demonstrates the use of functions, string manipulation, slicing, and concatenation in Python. It helps in understanding how strings can be dynamically modified using indexing and basic operations.

This exercise improved knowledge of string handling techniques and strengthened logical thinking in manipulating textual data using programming.

Overall, the objective of inserting a string into the middle of another string was successfully achieved using Python programming, highlighting the flexibility of string operations in Python.
}
\newcommand{\consixteen}{
In this assignment, a Python program was developed to sort the characters of a given string in lexicographical (dictionary) order. The program takes a string input from the user and processes it using built-in sorting techniques.

The logic of the program is based on breaking the string into individual characters, arranging them in ascending order based on their ASCII values, and then joining them back to form the sorted string.

The implementation demonstrates the use of string manipulation, sorting functions, and list operations in Python. It shows how built-in methods can be effectively used to simplify computational tasks.

This exercise helped in understanding how sorting works at a basic level and improved knowledge of string processing techniques in programming.

Overall, the objective of sorting a string lexicographically was successfully achieved using Python programming, highlighting the efficiency of Python’s built-in sorting capabilities.
}
\newcommand{\conseventeen}{
In this assignment, a Python program was developed to replace a substring within a given string without using built-in replace methods. The program takes three inputs: the main string, the string to be replaced, and the replacement string.

The logic is implemented by manually traversing the main string character by character and checking for a match with the target substring. When a match is found, it is replaced with the new string; otherwise, the characters are copied as they are.

The implementation demonstrates the use of loops, conditional statements, string traversal, and manual pattern matching in Python. It helps in understanding how string replacement works internally without relying on built-in functions.

This exercise improved logical thinking and enhanced understanding of string manipulation and algorithmic problem-solving techniques.

Overall, the objective of replacing a string without using built-in methods was successfully achieved using Python programming, highlighting the importance of manual algorithm design.
}
\newcommand{\coneighteen}{
In this assignment, a Python program was developed to concatenate two strings into a third string without using the + operator. The program takes two strings as input from the user and combines them using alternative string handling techniques.

The logic is implemented by iterating through each character of the first string and appending it to a result string, followed by iterating through each character of the second string and appending it to the same result variable. In this way, both strings are joined without directly using the concatenation operator.

The implementation demonstrates the use of loops, string traversal, and character-by-character processing in Python. It helps in understanding how string concatenation works at a fundamental level.
}
\newcommand{\connineteen}{
In this assignment, a Python program was developed to remove a set of specified characters from a given string. The program takes a string and a set of characters as input from the user and processes the string accordingly.

The logic is implemented by iterating through each character of the input string and checking whether it belongs to the set of characters to be removed. 
}
\newcommand{\contwenty}{
In this assignment, a Python program was developed to extract the first n characters from a given string. The program takes a string and an integer value n as input from the user.

The logic is implemented by traversing the string and selecting only the characters from index 0 to n-1. These characters are then stored in a new string and displayed as output.

The implementation demonstrates the use of string indexing, loops, and basic conditional logic in Python. It helps in understanding how parts of a string can be accessed and manipulated efficiently.

This exercise improved knowledge of string slicing concepts and strengthened basic programming skills in handling textual data.

Overall, the objective of extracting the first n characters of a string was successfully achieved using Python programming, highlighting the simplicity and effectiveness of string operations.
}

\newcommand{\contwentyone}{
In this assignment, a Python program was developed to generate a list of 10 random integers and separate them into two different lists: oddlist and evenlist. The program uses the random module to create random integer values.

After generating the list, each element is checked using a loop to determine whether it is odd or even. If a number is divisible by 2, it is added to the evenlist; otherwise, it is added to the oddlist.

The implementation demonstrates the use of lists, loops, conditional statements, and random number generation in Python.
}

\newcommand{\contwentytwo}{
In this assignment, a Python program was developed to sort a list of elements in ascending order using the insertion sort algorithm. The program repeatedly selects each element from the unsorted portion and places it into its correct position in the sorted portion of the list.

The implementation demonstrates the use of loops, nested iteration, and step-by-step comparison logic in Python. It helps in understanding how sorting algorithms work internally by building the sorted list incrementally.
}
\newcommand{\contwentythree}{
In this assignment, a Python program was developed to search for a key element in a sorted list using the binary search algorithm. The program repeatedly divides the search interval into two halves and compares the middle element with the target key to determine the next search direction.

The implementation demonstrates the use of conditional statements, loops, and efficient search techniques in Python. It highlights how binary search reduces the number of comparisons by eliminating half of the search space in each step.

This exercise improved understanding of algorithmic efficiency and strengthened knowledge of searching techniques, especially the importance of working with sorted data structures.

Overall, the objective of finding a key element using binary search was successfully achieved using Python programming, showcasing the efficiency of divide-and-conquer strategies.
}
\newcommand{\contwentyfour}{
In this assignment, a Python program was developed to create a list containing the first eight letters of the English alphabet and perform various slicing operations on it. The program demonstrates how list slicing can be used to extract specific portions of a list efficiently.

Using slicing techniques, the program prints the first three letters, selects any three letters from the middle of the list, and also displays elements from a given index up to the end of the list.

The implementation demonstrates the use of lists, indexing, and slicing operations in Python. It helps in understanding how data can be accessed and manipulated in different segments without modifying the original list.
}
\newcommand{\contwentyfive}{
In this assignment, a Python program was developed to sort a list of elements in ascending order using the selection sort algorithm. The program repeatedly selects the smallest element from the unsorted portion of the list and swaps it with the first unsorted element.

The implementation demonstrates the use of loops, nested iteration, swapping techniques, and comparison operations in Python. It helps in understanding how selection sort works by progressively building a sorted list from left to right.

This exercise improved problem-solving skills and strengthened the understanding of basic sorting algorithms and their step-by-step execution.

Overall, the objective of sorting elements in ascending order using selection sort was successfully achieved using Python programming, highlighting the importance of fundamental algorithm design.
}
\newcommand{\contwentysix}{
In this assignment, a Python program was developed to find the maximum value from the second half of a given list. The program divides the list into two equal parts and focuses only on the elements present in the second half.

The implementation demonstrates the use of list slicing, loops, and comparison operations in Python. It efficiently traverses the selected portion of the list to determine the largest element.

This exercise improved understanding of list manipulation and strengthened logical thinking in handling partial data structures and performing targeted computations.

Overall, the objective of finding the maximum value of the second half of the list was successfully achieved using Python programming, highlighting the effectiveness of list slicing and basic algorithmic operations.
}
\newcommand{\contwentyseven}{
In this assignment, a Python program was developed to generate a list of numbers from 1 to 100 that are divisible by either 5 or 6. The program iterates through the given range and checks each number using conditional statements.

If a number is divisible by 5 or 6, it is added to a result list. This approach demonstrates the use of loops, conditional logic, and list operations in Python.
}
\newcommand{\contwentyeight}{
In this assignment, a Python program was developed to accept five numbers from the user and compute their Greatest Common Divisor (GCD) using a function. The program takes user input and applies a systematic approach to determine the GCD of all the given numbers.

The implementation demonstrates the use of functions, loops, and mathematical logic in Python. It typically uses the Euclidean algorithm to compute the GCD efficiently by repeatedly finding the remainder until the final result is obtained.

This exercise helped in understanding how mathematical concepts like GCD can be implemented programmatically and reinforced the importance of modular programming using functions.
}
\newcommand{\contwentynine}{
In this assignment, a Python function named addfruit was developed to store fruit names along with their prices in a dictionary. The function takes a set of fruit names and their corresponding prices as input and returns a dictionary containing the entered information.

The program also includes exception handling using the ValueError exception. If a fruit name already exists in the dictionary, the program raises a ValueError to prevent duplicate entries.

The implementation demonstrates the use of functions, dictionaries, condition checking, and exception handling in Python. It helps in understanding how data integrity can be maintained using error handling techniques.

This exercise improved knowledge of dictionary operations and strengthened the understanding of how exceptions can be used to manage invalid or duplicate inputs in a program.

Overall, the objective of creating a function with duplicate checking and exception handling was successfully achieved using Python programming, highlighting the importance of robust program design.
}
\newcommand{\conthirty}{
In this assignment, a Python function was developed to store Air Quality Index (AQI) values with the corresponding dates as keys in a dictionary. The program takes user input for the date and AQI value and stores them in a structured key-value format.

The implementation demonstrates the use of functions, dictionaries, and user input handling in Python. It shows how real-world environmental data can be organized efficiently using key-value pairs.
}

\newcommand{\conthirtyone}{
In this assignment, a Python program was developed to create a dictionary containing usernames as keys and passwords as values. The program initializes five sample entries to demonstrate basic dictionary operations in Python.

The implementation also demonstrates the use of the clear() method, which removes all items from a dictionary, and the fromkeys() method, which creates a new dictionary with specified keys and a common default value.
}
\newcommand{\conthirtytwo}{
In this assignment, a Python program was developed to create a dictionary where country names are used as keys and their corresponding capital cities are stored as values. The program takes user input and stores multiple entries in a dictionary structure.

The implementation then displays the dictionary contents in sorted order based on the country names. This is achieved by applying sorting techniques on the dictionary keys before printing the results.
}
\newcommand{\conthirtythree}{
In this assignment, a Python program was developed to create a dictionary containing friends’ names as keys and their phone numbers as values. The program stores multiple entries and displays the dictionary in sorted order based on the names.

The implementation also allows the user to search for a specific name in the dictionary. If the name is found, the corresponding phone number is displayed; otherwise, the user is prompted to enter the new contact details, which are then added to the dictionary.

The program demonstrates the use of dictionaries, sorting, conditional statements, and user input handling in Python. It helps in understanding how real-world contact management systems can be implemented using simple data structures.

This exercise improved knowledge of dictionary operations and strengthened logical thinking in handling dynamic data insertion and retrieval.

Overall, the objective of creating and managing a friends’ contact dictionary with search and insertion functionality was successfully achieved using Python programming.
}
\newcommand{\conthirtyfour}{
In this assignment, a Python program was developed to create a dictionary containing author names as keys and their corresponding ISBN numbers as values. The program includes five sample entries to demonstrate basic dictionary creation and manipulation.

The implementation also demonstrates the use of the clear() method, which removes all entries from the dictionary, and the fromkeys() method, which creates a new dictionary using specified keys with a common default value.
}
\newcommand{\conthirtyfive}{
In this assignment, a Python program was developed to create a list containing three elements and then convert that list into a tuple. The program demonstrates the basic concept of type conversion between mutable and immutable data structures in Python.

The implementation shows how a list, which is mutable, can be easily transformed into a tuple, which is immutable, using built-in Python functions. This helps in understanding the difference between these two data structures.

The exercise demonstrates the use of lists, tuples, and type conversion methods in Python. It also improves understanding of how data structures can be manipulated and converted based on program requirements.
}
\newcommand{\conthirtysix}{
In this assignment, a Python program was developed to demonstrate tuple unpacking by assigning the elements of a tuple to multiple variables. The program creates a tuple containing several values and then extracts each value into separate variables.

The implementation shows how Python supports multiple assignment, where each element of a tuple can be directly mapped to individual variables in a single statement. This eliminates the need for manual indexing and improves code readability.

The exercise demonstrates the use of tuples, variable assignment, and data unpacking techniques in Python. It helps in understanding how structured data can be efficiently distributed into separate variables.

This exercise improved understanding of tuple operations and strengthened knowledge of Python’s built-in unpacking feature.

Overall, the objective of unpacking a tuple into multiple variables was successfully achieved using Python programming, highlighting the simplicity and efficiency of this feature.
}
\newcommand{\conthirtyseven}{
In this assignment, a Python program was developed to check whether a given item exists within a tuple. The program takes a tuple of elements and an item to be searched as input, and then verifies the presence of the item.

The implementation uses membership operators in Python to determine whether the specified element is present in the tuple or not. Based on the result, an appropriate message is displayed to the user.

The exercise demonstrates the use of tuples, membership testing, and conditional statements in Python. It helps in understanding how data can be efficiently searched within immutable data structures.

This exercise improved logical thinking and strengthened knowledge of tuple operations and searching techniques in programming.

Overall, the objective of checking whether an item exists in a tuple was successfully achieved using Python programming, highlighting the simplicity and effectiveness of membership operations.
}
\newcommand{\conthirtyeight}{
In this assignment, a Python program was developed to unzip a list of tuples into individual lists. The program takes a list containing multiple tuples and separates each position of the tuples into separate lists.

The implementation demonstrates the use of tuple unpacking and built-in Python techniques such as the zip(*) function. This allows the data to be efficiently transformed from a grouped structure into separate collections.

The exercise demonstrates the use of tuples, lists, and data unpacking concepts in Python. It helps in understanding how structured data can be reorganized using simple built-in operations.

This exercise improved knowledge of data manipulation techniques and strengthened understanding of how to handle and transform complex data structures.

Overall, the objective of unzipping a list of tuples into individual lists was successfully achieved using Python programming, highlighting the efficiency of Python’s built-in functions.
}
\newcommand{\conthirtynine}{
In this assignment, a Python program was developed to perform basic set operations such as intersection, union, difference, and symmetric difference between two sets. The program takes two sets as input and applies standard set operations to obtain the required results.

The implementation demonstrates the use of Python sets and built-in operators such as union ($|$), intersection ($&$), difference ($-$), and symmetric difference. These operations help in efficiently comparing and combining data elements.

The exercise demonstrates the use of sets, mathematical operations, and logical reasoning in Python. It helps in understanding how different relationships between data collections can be analyzed using set theory concepts.
}
\newcommand{\conforty}{
In this assignment, a Python program was developed to demonstrate the use of issubset() and issuperset() methods in sets. The program defines two sets and checks their relationship using these built-in methods.

The issubset() method is used to determine whether all elements of one set are present in another set, while issuperset() checks whether a set contains all elements of another set. Based on these checks, the program displays appropriate results.

The implementation demonstrates the use of sets, relational operations, and built-in set methods in Python. It helps in understanding how relationships between different sets can be evaluated efficiently.
}

\newcommand{\confortyone}{
In this assignment, a Python program was developed to take a range of numbers and create a list of tuples where each tuple contains a number and its square. The program generates values within the given range and processes each number sequentially.

The implementation demonstrates the use of loops, tuple creation, and arithmetic operations in Python. Each number in the range is paired with its square using exponentiation and stored in a list of tuples.

The exercise demonstrates the use of lists, tuples, and mathematical operations in Python. It helps in understanding how structured data can be generated and stored efficiently using iteration.

This exercise improved logical thinking and strengthened knowledge of tuple handling and range-based computations in programming.

Overall, the objective of creating a list of tuples containing numbers and their squares was successfully achieved using Python programming, highlighting the effectiveness of iterative data generation.
}

\newcommand{\confortytwo}{
In this assignment, a Python program was developed to demonstrate how to clear a set using built-in methods. The program first creates a set containing multiple elements and then removes all the elements using the clear() method.

The implementation shows how the clear() function empties the entire set while preserving the set object itself. After performing the operation, the program verifies that the set has become empty.

The exercise demonstrates the use of sets and built-in methods in Python. It helps in understanding how data structures can be reset or cleared efficiently without deleting the variable itself.
}

\newcommand{\confortythree}{
In this assignment, a Python program was developed to determine the length of a set, which represents the total number of unique elements present in it. The program creates a set with multiple elements and then calculates its size.

The implementation uses the built-in len() function to find the number of elements stored in the set. Since sets do not allow duplicate values, only unique items are counted.

The exercise demonstrates the use of sets and built-in functions in Python. It helps in understanding how data structures can be analyzed to obtain useful information such as size or cardinality.

This exercise improved knowledge of set operations and strengthened understanding of basic data structure properties in Python.
}

\newcommand{\confortyfour}{
In this assignment, a Python program was developed to store the latitude and longitude of a location using a tuple and display the stored values. The program defines a tuple containing two floating-point values representing geographical coordinates.

The implementation demonstrates the use of tuples in Python to store immutable data. It shows how multiple related values can be grouped together in a single data structure and accessed easily when required.

The exercise demonstrates the use of tuples, data storage, and output operations in Python. It helps in understanding how real-world information such as geographical coordinates can be represented in programming.
}

\newcommand{\confortyfive}{
In this assignment, a Python program was developed to read a text file provided by the user, extract all the words, and display the non-duplicate words in ascending order. The program processes the file content by splitting it into individual words.

The implementation uses sets to remove duplicate words and sorting techniques to arrange the unique words in alphabetical order. It demonstrates how file handling and data processing can be combined effectively in Python.

The exercise demonstrates the use of file handling, sets, loops, and sorting operations in Python. It helps in understanding how textual data can be processed and organized efficiently.
}

\newcommand{\confortysix}{
In this assignment, a Python program was developed to determine the size of a plain text file. The program takes the file as input and calculates its size in bytes.

The implementation demonstrates the use of file handling techniques in Python. It accesses the file properties using built-in methods or modules such as os, which allow retrieval of file size information efficiently.

The exercise demonstrates the use of file operations and system-level functions in Python. It helps in understanding how file metadata can be accessed and utilized in programming.
}

\newcommand{\confortyseven}{
In this assignment, a Python program was developed to read a text file provided by the user and count the number of vowels and consonants present in its content. The program processes the file by reading its text and analyzing each character individually.

The implementation uses file handling, loops, and conditional statements in Python. Each character is checked to determine whether it is a vowel or consonant, and the corresponding counters are updated accordingly.
}

\newcommand{\confortyeight}{
In this assignment, a Python program was developed to read the first n lines of a text file based on user input. The program prompts the user to enter the value of n and then reads the specified number of lines from the file.

The implementation demonstrates the use of file handling, loops, and user input in Python. It reads the file line by line and displays only the required number of lines as specified by the user.

The exercise demonstrates how file reading operations can be controlled and customized using simple logic in Python. It helps in understanding efficient file processing techniques without loading the entire file into memory.

This exercise improved knowledge of file handling and strengthened understanding of controlled data reading and basic input-output operations.

Overall, the objective of reading the first n lines of a file was successfully achieved using Python programming, highlighting its usefulness in efficient file processing.
}

\newcommand{\confortynine}{
In this assignment, a Python program was developed to read the contents of a text file and count the occurrences of each letter present in it. The program processes the file by reading its text and analyzing each character individually.

The implementation uses file handling, loops, and dictionary data structures in Python. Each letter is checked and its frequency is stored in a dictionary, which helps in maintaining a count of all occurrences efficiently.

The exercise demonstrates the use of file operations, string processing, and data analysis using dictionaries in Python. It helps in understanding how textual data can be analyzed and summarized programmatically.

This exercise improved knowledge of file handling and strengthened understanding of frequency analysis and dictionary usage in Python programming.

Overall, the objective of counting the occurrences of each letter in a file was successfully achieved using Python programming, highlighting its usefulness in text analysis tasks.
}

\newcommand{\confifty}{
In this assignment, a Python program was developed to read the last n lines of a text file based on user input. The program prompts the user to enter the value of n and then retrieves the required number of lines from the end of the file.

The implementation demonstrates the use of file handling, list operations, and user input in Python. The file content is read and processed in such a way that only the last n lines are selected and displayed.

The exercise demonstrates efficient file processing techniques and shows how data can be accessed from the end of a file without manually iterating through the entire content line by line.

This exercise improved understanding of file handling and strengthened knowledge of list manipulation and data retrieval methods in Python.

Overall, the objective of reading the last n lines of a file was successfully achieved using Python programming, highlighting its efficiency in handling file-based operations.
}

\newcommand{\confiftyone}{
In this assignment, a Python program was developed to combine each line from the first file with the corresponding line in the second file. The program reads both files simultaneously and processes them line by line.

The implementation uses file handling techniques and iteration in Python. Each line from the first file is paired with the matching line from the second file, and the combined result is displayed or stored accordingly.

The exercise demonstrates the use of file operations, looping, and parallel processing of data from multiple sources in Python. It helps in understanding how data from different files can be merged efficiently.
}

\newcommand{\confiftytwo}{
In this assignment, a Python program was developed to remove newline characters from a text file. The program reads the file content and processes it to eliminate all line breaks.

The implementation demonstrates the use of file handling, string manipulation, and replacement techniques in Python. By replacing newline characters with spaces or an empty string, the text is converted into a continuous format.

The exercise demonstrates how file content can be cleaned and formatted using simple string operations. It helps in understanding basic text preprocessing techniques in programming.
}

\newcommand{\confiftythree}{
In this assignment, a Python program was developed to read a random line from a text file. The program loads the file content and selects one line randomly for display.

The implementation demonstrates the use of file handling and random module functions in Python. The file is read line by line, and a random index is generated to pick a line from the list of lines.

The exercise demonstrates how file operations can be combined with randomness to retrieve unpredictable data from a dataset. It helps in understanding basic file processing and random selection techniques in programming.

This exercise improved knowledge of file handling and strengthened understanding of list operations and randomization concepts in Python.

Overall, the objective of reading a random line from a file was successfully achieved using Python programming, highlighting its usefulness in file-based data selection tasks.
}
\newcommand{\confiftyfour}{
In this assignment, a Python program was developed to read data from one CSV file and write the same contents into another CSV file. The program uses file handling techniques to access and process structured tabular data.

The implementation demonstrates the use of the csv module in Python, which allows efficient reading and writing of CSV files. The data is read row by row from the source file and then written into the destination file without modification.

The exercise demonstrates the use of file handling, CSV processing, and data transfer techniques in Python. It helps in understanding how structured data can be copied and managed between files.
}
\newcommand{\confiftyfive}{
In this assignment, a Python program was developed to identify and match words that contain the letter 'z'. The program processes a given input text or list of words and checks each word for the presence of the specified character.

The implementation demonstrates the use of string handling and pattern matching techniques in Python. It typically uses regular expressions or simple conditional checks to determine whether a word contains the letter 'z'.

The exercise demonstrates the use of pattern matching, string processing, and basic text analysis in Python. It helps in understanding how specific patterns can be identified within textual data.
}
\newcommand{\confiftysix}{
In this assignment, a Python program was developed to remove all leading zeros from an IP address. The program takes an IP address as input and processes each octet individually.

The implementation splits the IP address into its components using the dot (.) separator, converts each part into an integer to automatically remove leading zeros, and then reconstructs the cleaned IP address in proper format.

The exercise demonstrates the use of string manipulation, type conversion, and list operations in Python. It helps in understanding how formatted numeric data can be cleaned and standardized.

This exercise improved knowledge of string processing and strengthened understanding of data cleaning techniques in real-world applications such as networking.

Overall, the objective of removing leading zeros from an IP address was successfully achieved using Python programming, highlighting its usefulness in data formatting and validation tasks.
}
\newcommand{\confiftyseven}{
In this assignment, a Python program was developed to search for numbers ($0–9$) of length between 1 to 3 digits in a given string. The program processes the input text and identifies numeric patterns that satisfy the given condition.

The implementation demonstrates the use of regular expressions in Python for pattern matching. It efficiently searches for sequences of digits ranging from single-digit numbers to three-digit numbers within the input string.
}
\newcommand{\confiftyeight}{
In this assignment, a Python program was developed to find substrings within a given string. The program takes a main string and a target substring as input and checks whether the substring is present.

The implementation demonstrates the use of string handling techniques in Python. It uses basic searching logic or built-in string methods to identify the presence and position of substrings within the main string.

The exercise demonstrates the use of strings, searching operations, and text processing in Python. It helps in understanding how smaller patterns can be located inside larger text data.

}
\newcommand{\confiftynine}{
In this assignment, a Python program was developed to extract the year, month, and date from a URL. The program processes the URL string and identifies the date pattern embedded within it.

The implementation demonstrates the use of string manipulation and regular expressions in Python. It searches for a structured date format (such as YYYY-MM-DD) and separates it into year, month, and day components.

The exercise demonstrates the use of pattern matching, string processing, and data extraction techniques in Python. It helps in understanding how structured information can be retrieved from unstructured or semi-structured text.
}
\newcommand{\consixty}{
In this assignment, a Python program was developed to read a text file and convert dates from the format yyyy-mm-dd to dd-mm-yyyy. The program processes the file content and identifies date patterns for transformation.

The implementation demonstrates the use of file handling and string manipulation techniques in Python. It uses splitting or regular expressions to extract the year, month, and day components and rearranges them into the required format.

The exercise demonstrates the use of text processing, pattern recognition, and data transformation in Python. It helps in understanding how structured data formats can be modified programmatically.
}

\newcommand{\consixtyone}{
In this assignment, a Python program was developed to abbreviate the word "Street" as "St." in a given string. The program takes an input string from the user and processes it to replace occurrences of the target word.

The implementation demonstrates the use of string manipulation techniques in Python. It identifies the word "Street" within the text and replaces it with its abbreviated form "St." using basic string processing or pattern matching methods.
}
\newcommand{\consixtytwo}{
In this assignment, a Python program was developed to find all words that are exactly five characters long in a given string. The program takes a sentence or text input from the user and processes it to identify matching words.

The implementation demonstrates the use of string manipulation and pattern matching techniques in Python. It splits the text into individual words and checks the length of each word to filter out those with exactly five characters.

}
\newcommand{\consixtythree}{
In this assignment, a Python class named Quadratic was developed to represent a quadratic equation with data attributes a, b, and c. The class includes an \_\_init\_\_() method to initialize these coefficients and a roots() method to compute the roots of the quadratic equation.

The roots are calculated using the standard quadratic formula:

\[
x = \frac{-b \pm \sqrt{b^2 - 4ac}}{2a}
\]

The implementation demonstrates the use of object-oriented programming concepts in Python, including class creation, constructors, and methods. It also makes use of mathematical operations to determine real or complex roots based on the discriminant.
}
\newcommand{\consixtyfour}{
In this assignment, a Python program was developed using object-oriented programming concepts to define a class named Student and display the marks details of the top five students using inheritance. The base class is used to store basic student information and marks, while the derived class extends its functionality to process and display results.

The implementation demonstrates the use of classes, inheritance, methods, and data handling in Python. It also includes logic to determine and display the top five students based on their marks.

The exercise demonstrates how object-oriented principles such as inheritance can be applied to real-world scenarios like student data management. It helps in understanding code reusability and structured program design.

This exercise improved knowledge of OOP concepts and strengthened understanding of inheritance and data organization in Python programming.

Overall, the objective of defining a student class and displaying top five student details using inheritance was successfully achieved using Python programming, highlighting the power of object-oriented design.
}
\newcommand{\consixtyfive}{
In this assignment, a Python program was developed using object-oriented programming to create a class named Library with data attributes such as acc\_number, title, author, and publisher. The class includes methods to handle book details and fine calculation.

The read() method is used to accept and store the book information including accession number, title, and author. The compute() method calculates the fine based on the number of days a book is returned late, using the rate of $1.50 per day.

The implementation demonstrates the use of classes, methods, data attributes, and basic arithmetic operations in Python. It shows how real-world systems like library management can be modeled using object-oriented concepts.

This exercise improved understanding of class design and strengthened knowledge of applying OOP principles to practical problem-solving.

Overall, the objective of creating a library class and calculating late return fines was successfully achieved using Python programming, highlighting the effectiveness of object-oriented design in real-life applications.
}
\newcommand{\consixtysix}{
In this assignment, a Python program was developed using multiple inheritance to create a class named CalendarClock, which inherits properties from two base classes: Clock and Calendar. The Clock class is responsible for handling time-related information, while the Calendar class manages date and month details.

The derived class CalendarClock combines the features of both parent classes and displays complete information including date, month, and current time in a structured format.

The implementation demonstrates the use of classes, multiple inheritance, methods, and object-oriented programming principles in Python. It helps in understanding how multiple functionalities can be combined efficiently into a single class.

This exercise improved knowledge of object-oriented programming concepts such as inheritance hierarchy and code reusability, and strengthened the ability to design modular programs.

Overall, the objective of creating a CalendarClock class using multiple inheritance to display date and time was successfully achieved using Python programming, highlighting the power of OOP in real-world applications.
}
\newcommand{\consixtyseven}{
In this assignment, a Python program was developed using object-oriented programming to add two polynomials using classes. Each polynomial is represented as an object, and the coefficients are stored as data attributes.

The program defines methods to accept polynomial terms and perform addition by combining corresponding coefficients of the two polynomials. The result is then displayed as a new polynomial expression.

The implementation demonstrates the use of classes, objects, lists, and arithmetic operations in Python. It helps in understanding how mathematical expressions like polynomials can be represented and manipulated using programming concepts.
}
\newcommand{\consixtyeight}{
In this assignment, JSON (JavaScript Object Notation) was studied as a lightweight data-interchange format used for storing and exchanging structured data. It is widely used in web applications and APIs due to its simplicity and readability.

A simple JSON document was constructed containing key-value pairs representing sample data such as name, age, and city. The Python program then parsed this JSON document using the built-in json module.

The implementation demonstrates the use of JSON handling, parsing techniques, and built-in library functions in Python. It helps in understanding how structured data can be converted between JSON format and Python objects such as dictionaries.

This exercise improved knowledge of data formats and strengthened understanding of data exchange mechanisms in programming.

Overall, the objective of defining JSON, constructing a JSON document, and parsing it using Python was successfully achieved, highlighting its importance in modern data communication.
}
\newcommand{\consixtynine}{
In this assignment, the differences between XML (Extensible Markup Language) and JSON (JavaScript Object Notation) were studied and analyzed. Both formats are widely used for data storage and exchange in software applications.

XML is a markup language that uses tags to define data structure, making it more verbose and self-descriptive. It supports attributes and nested elements, which makes it suitable for complex document representation. However, it is relatively heavier in size and more complex to parse.

On the other hand, JSON is a lightweight data-interchange format that uses key-value pairs. It is easier to read, write, and parse, especially in programming languages like Python and JavaScript. JSON is more compact compared to XML and is commonly used in web APIs and modern applications.
}

\newcommand{\conseventy}{
In this assignment, XML (Extensible Markup Language) was studied as a markup language used to store and transport data in a structured and hierarchical format. XML is widely used in data exchange systems due to its flexibility and self-descriptive nature.

A simple XML document was constructed containing sample data with nested elements to represent structured information. The Python program then processed this XML file using the ElementTree module.
}

\newcommand{\conseventyone}{
In this assignment, various NumPy array creation functions were studied and implemented using Python. Functions such as array(), zeros(), ones(), arange(), linspace(), and random() were used to create arrays with different structures and values.
}
\newcommand{\conseventytwo}{
In this assignment, various NumPy indexing techniques including integer indexing, array indexing, Boolean array indexing, and slicing were studied and implemented using Python. These techniques provide powerful ways to access, modify, and manipulate elements within NumPy arrays.

The implementation demonstrated how integer indexing can be used to access specific elements, array indexing can retrieve multiple elements simultaneously, Boolean indexing can filter data based on conditions, and slicing can extract subsets of arrays efficiently.

The exercise highlighted the flexibility and efficiency of NumPy in handling large datasets and performing advanced data manipulation operations. It also helped in understanding how different indexing methods can be applied to solve practical computational problems.

}
\newcommand{\conseventythree}{
In this assignment, a Python program was developed to create and display a one-dimensional array-like object using the Pandas library. The program utilizes the Series data structure, which is one of the fundamental components of Pandas for handling labeled data.
}
\newcommand{\conseventyfour}{
In this assignment, a Python program was developed to perform arithmetic operations such as addition, subtraction, multiplication, and division on two Pandas Series. The program creates two Series objects and applies the respective operations element-wise.

The implementation demonstrates how Pandas automatically aligns data based on index values and performs arithmetic computations efficiently. It highlights the simplicity of carrying out mathematical operations on structured data using Pandas.
}

\newcommand{\conseventyfive}{
In this assignment, a Python program was developed to create and display a Pandas DataFrame from dictionary data with specified index labels. The program organizes data stored in a dictionary into a tabular structure using the DataFrame class.

The implementation demonstrates how dictionary keys become column names and how custom index labels can be assigned to rows for better identification and organization of data. The resulting DataFrame provides a clear and structured representation of the dataset.

The exercise demonstrates the use of the Pandas library, dictionaries, DataFrames, and index labels in Python. It helps in understanding how real-world data can be efficiently stored, organized, and displayed in tabular form.

This exercise improved knowledge of Pandas data structures and strengthened understanding of data representation and manipulation techniques in Python.

Overall, the objective of creating and displaying a DataFrame from dictionary data with index labels was successfully achieved, highlighting the effectiveness of Pandas for data analysis and management tasks.
}
\newcommand{\conseventysix}{
In this assignment, a Python program was developed to find the eigenvalues of a 3×3 matrix. The program accepts or defines a matrix and applies numerical computation techniques to determine its eigenvalues.

The implementation demonstrates the use of matrix operations and linear algebra concepts in Python, typically through the NumPy library. Eigenvalues are calculated by solving the characteristic equation of the matrix, which provides important information about its properties and behavior.

The exercise demonstrates the application of linear algebra in programming and helps in understanding how matrices can be analyzed using computational tools. Eigenvalues play a significant role in various fields such as machine learning, computer graphics, physics, and engineering.

This exercise improved knowledge of matrix operations and strengthened understanding of numerical computation and linear algebra concepts in Python.

Overall, the objective of finding the eigenvalues of a 3×3 matrix was successfully achieved using Python programming, highlighting the effectiveness of computational methods in solving mathematical problems.
}
\newcommand{\conseventyseven}{
In this assignment, a NumPy program was developed to compute the dot product of two vectors. The program creates two vectors and performs the dot product operation using NumPy's built-in functions.

The implementation demonstrates the use of vector operations and numerical computation techniques in Python. The dot product is calculated by multiplying the corresponding elements of the vectors and summing the resulting products.

The exercise highlights the importance of vector mathematics in scientific computing, machine learning, data analysis, and engineering applications. It also demonstrates how NumPy simplifies complex mathematical operations through efficient array handling.

This exercise improved knowledge of NumPy arrays and strengthened understanding of vector operations and linear algebra concepts in Python.

Overall, the objective of computing the dot product of two vectors was successfully achieved using NumPy, highlighting its effectiveness and efficiency for numerical and mathematical computations.
}
\newcommand{\conseventyeight}{
In this assignment, a Python program was developed to remove duplicate rows from a Pandas DataFrame based on a specific column. The program identifies repeated values in the selected column and retains only the required unique records.

The implementation demonstrates the use of Pandas DataFrame operations and data-cleaning techniques. By applying the appropriate method, duplicate entries are efficiently detected and removed while preserving the integrity of the dataset.

The exercise highlights the importance of data preprocessing in data analysis and machine learning tasks. It helps in understanding how duplicate records can affect the quality of analysis and how they can be handled effectively.

This exercise improved knowledge of Pandas DataFrame manipulation and strengthened understanding of data cleaning and preparation techniques in Python.
}

\newcommand{\conseventynine}{
In this assignment, a NumPy program was developed to create a 10×10 checkerboard matrix. The program generates a two-dimensional array in which adjacent elements alternate between 0 and 1, forming a checkerboard-like pattern.

The implementation demonstrates the use of NumPy array creation, indexing, and slicing techniques to efficiently construct the matrix. By utilizing NumPy's powerful array operations, the checkerboard pattern is generated with minimal code and high computational efficiency.

The exercise highlights the importance of matrix manipulation and pattern generation in numerical computing. It also provides practical experience in working with multidimensional arrays and array indexing concepts.
}
\newcommand{\coneighty}{
In this assignment, a Python program was developed to create a scatter plot with a regression line using the Altair visualization library. The program plots the relationship between two variables and overlays a regression line to illustrate the overall trend in the data.

The implementation demonstrates the use of Altair for interactive data visualization and statistical transformation. The scatter plot displays individual data points, while the regression line provides a visual representation of the linear relationship between the variables.

The exercise highlights the importance of data visualization in understanding patterns, trends, and correlations within datasets. It also demonstrates how Altair simplifies the process of creating informative and visually appealing charts.

This exercise improved knowledge of data visualization techniques and strengthened understanding of regression analysis and graphical data representation in Python.

Overall, the objective of plotting a scatter plot with a regression line using Altair was successfully achieved, highlighting the effectiveness of visualization tools in exploratory data analysis.
}

\newcommand{\coneightyone}{
In this assignment, a NumPy program was developed to compute the covariance matrix of a dataset. The program processes the dataset and calculates the covariance values that describe how different variables vary with respect to one another.

The implementation demonstrates the use of NumPy's statistical functions and matrix operations to efficiently compute the covariance matrix. The resulting matrix provides valuable information about the relationships and dependencies between variables in the dataset.

The exercise highlights the importance of covariance analysis in statistics, data science, and machine learning. It helps in understanding the direction and strength of relationships between variables and serves as a foundation for advanced techniques such as Principal Component Analysis (PCA).

This exercise improved knowledge of statistical computation and strengthened understanding of matrix-based data analysis using NumPy.

Overall, the objective of computing the covariance matrix of a dataset was successfully achieved using NumPy, demonstrating its effectiveness for statistical and numerical analysis tasks.
}
\newcommand{\coneightytwo}{
In this assignment, a Python program was developed to group a Pandas DataFrame by two columns and compute aggregate statistics. The program organizes the data into groups based on the unique combinations of values in the selected columns.

The implementation demonstrates the use of the groupby() function along with aggregate operations such as sum, mean, count, minimum, and maximum. These operations help summarize and analyze large datasets efficiently.

The exercise highlights the importance of data aggregation and grouping techniques in data analysis. It helps in understanding how meaningful insights can be extracted by categorizing data and computing statistical measures for each group.

This exercise improved knowledge of Pandas DataFrame operations and strengthened understanding of data summarization, grouping, and statistical analysis techniques in Python.

Overall, the objective of grouping a DataFrame by two columns and computing aggregate statistics was successfully achieved using Pandas, demonstrating its effectiveness for data analysis and reporting tasks.
}
\newcommand{\coneightythree}{
In this assignment, a Python program was developed to create a cumulative sum column in a Pandas DataFrame. The program processes a numerical column and calculates the running total of its values, storing the results in a new column.

The implementation demonstrates the use of Pandas DataFrame operations and the cumulative sum functionality. Each value in the new column represents the sum of all previous values up to and including the current row.

The exercise highlights the importance of cumulative calculations in data analysis, financial reporting, and time-series processing. It helps in understanding how trends and progressive totals can be tracked efficiently within a dataset.

This exercise improved knowledge of Pandas DataFrame manipulation and strengthened understanding of cumulative operations and data transformation techniques in Python.
}
\newcommand{\coneightyfour}{
In this assignment, a Python program was developed to generate 1000 random numbers from a normal distribution and visualize the data using a histogram created with the Altair library. The generated dataset follows the characteristics of a normal distribution, making it suitable for statistical analysis and visualization.
}
\newcommand{\coneightyfive}{
In this assignment, a Python function was developed to standardize a NumPy array so that the resulting data has a mean of 0 and a variance of 1. The function computes the mean and standard deviation of the array and then transforms each element using the standardization formula.

The implementation demonstrates the use of NumPy's statistical functions and array operations to perform data normalization efficiently. Standardization ensures that the data is scaled consistently, making it suitable for statistical analysis and machine learning applications.

The exercise highlights the importance of data preprocessing in numerical computing and predictive modeling. It helps in understanding how differences in scale can be removed while preserving the underlying distribution of the data.

This exercise improved knowledge of NumPy arrays, statistical measures, and data transformation techniques. It also strengthened understanding of feature scaling, which is an essential step in many data science and machine learning workflows.

Overall, the objective of standardizing a NumPy array to achieve a mean of 0 and a variance of 1 was successfully achieved using Python and NumPy, demonstrating the effectiveness of mathematical transformations in data preprocessing.
}
\newcommand{\coneightysix}{
In this assignment, a Python program was developed to reshape a one-dimensional NumPy array into a three-dimensional array of appropriate dimensions. The program first creates a 1D array and then reorganizes its elements into a 3D structure using NumPy's reshape() function.

The implementation demonstrates the use of array reshaping techniques and multidimensional data structures in NumPy. The reshape operation preserves all elements while changing the arrangement of data into multiple dimensions.

The exercise highlights the importance of array manipulation in numerical computing, data analysis, and machine learning. It helps in understanding how data can be transformed into different shapes to meet the requirements of various computational tasks.

This exercise improved knowledge of NumPy array operations and strengthened understanding of multidimensional arrays and data representation techniques in Python.

Overall, the objective of reshaping a one-dimensional NumPy array into a three-dimensional array was successfully achieved using NumPy, demonstrating its flexibility and efficiency in handling complex data structures.
}
\newcommand{\coneightyseven}{
In this assignment, a Python program was developed to merge three Pandas DataFrames using a common key. The program combines data from multiple sources into a single unified DataFrame based on matching key values.

The implementation demonstrates the use of Pandas merge operations, which efficiently join rows from different DataFrames while maintaining the relationships between corresponding records. This technique is widely used for integrating and analyzing data stored across multiple tables.

The exercise highlights the importance of data integration and relational data processing in data analysis. It helps in understanding how information from different datasets can be combined to create a more comprehensive view of the data.

This exercise improved knowledge of Pandas DataFrame operations and strengthened understanding of data merging, joining, and management techniques in Python.
}
\newcommand{\coneightyeight}{
In this assignment, a Python program was developed to create an interactive chart using the Altair visualization library with a dropdown selection filter. The chart allows users to dynamically select values from a dropdown menu and view the corresponding filtered data.

The implementation demonstrates the use of Altair's interactive features, including selection parameters, dropdown widgets, and dynamic data filtering. These features enhance user interaction and make data exploration more intuitive and efficient.

The exercise highlights the importance of interactive visualizations in data analysis and decision-making. It helps in understanding how user-controlled filters can be integrated into charts to provide deeper insights into datasets.

This exercise improved knowledge of Altair visualization techniques and strengthened understanding of interactive chart design and data filtering concepts in Python.

Overall, the objective of creating an interactive Altair chart with a dropdown selection filter was successfully achieved, demonstrating the effectiveness of interactive visualizations for exploratory data analysis.
}
\newcommand{\coneightynine}{
In this assignment, a Python program was developed to detect and remove outliers from a dataset using the Z-score method in Pandas. The program calculates the Z-score for each data point and identifies observations that lie significantly far from the mean.

The implementation demonstrates the use of statistical analysis and data preprocessing techniques in Python. Records with Z-score values beyond a specified threshold are considered outliers and are removed from the dataset to improve data quality.

The exercise highlights the importance of outlier detection in data analysis, machine learning, and statistical modeling. Removing extreme values helps reduce bias and improves the reliability of analytical results.
}
\newcommand{\conninety}{
In this assignment, a Python program was developed to compute the pairwise correlation matrix of a Pandas DataFrame. The program analyzes the relationships between numerical columns and calculates correlation coefficients for every possible pair of variables.

The implementation demonstrates the use of Pandas statistical functions to efficiently generate a correlation matrix. The resulting matrix provides a clear representation of the strength and direction of relationships among the variables in the dataset.

The exercise highlights the importance of correlation analysis in statistics, data science, and machine learning. It helps in identifying patterns, dependencies, and potential associations between variables, which can support better decision-making and predictive modeling.

This exercise improved knowledge of Pandas DataFrame operations and strengthened understanding of statistical analysis and data exploration techniques in Python.

Overall, the objective of computing the pairwise correlation matrix of a DataFrame was successfully achieved using Pandas, demonstrating its effectiveness in analyzing relationships within structured datasets.
}

\newcommand{\conninetyone}{
In this assignment, a Python program was developed to create a heatmap visualization using the Altair library. The heatmap represents data values through variations in color intensity, making it easier to identify patterns, trends, and relationships within the dataset.

The implementation demonstrates the use of Altair's visualization capabilities, including encoding data values into color scales and arranging them in a grid-like structure. This approach provides an intuitive and visually appealing way to analyze multidimensional data.

}
\newcommand{\conninetytwo}{
In this assignment, a Python program was developed to apply a custom function to each element of a NumPy array using vectorization. The program defines a user-created function and efficiently applies it across all elements of the array without using explicit loops.

The implementation demonstrates the use of NumPy's vectorization capabilities, which allow element-wise operations to be performed quickly and efficiently. By utilizing vectorized functions, the program achieves better performance and cleaner code compared to traditional iterative approaches.

The exercise highlights the importance of vectorized computation in numerical analysis, scientific computing, and data processing. It helps in understanding how large datasets can be manipulated efficiently using optimized array operations.

This exercise improved knowledge of NumPy arrays and strengthened understanding of vectorization techniques and function application in Python.

Overall, the objective of applying a custom function to each element of a NumPy array using vectorization was successfully achieved, demonstrating the power and efficiency of NumPy for high-performance numerical computations.
}
\newcommand{\conninetythree}{
In this assignment, a Python program was developed to convert a categorical column into one-hot encoded format using the Pandas library. The program transforms categorical values into multiple binary columns, where each category is represented by a separate column.

The implementation demonstrates the use of Pandas data preprocessing techniques, particularly the creation of dummy variables for categorical data. This transformation enables categorical information to be represented numerically while preserving the distinct categories.

The exercise highlights the importance of data preprocessing in machine learning and data analysis. One-hot encoding is commonly used to prepare categorical variables for algorithms that require numerical input.

This exercise improved knowledge of Pandas DataFrame operations and strengthened understanding of feature engineering and data transformation techniques in Python.

Overall, the objective of converting a categorical column into one-hot encoded format was successfully achieved using Pandas, demonstrating its effectiveness in preparing data for analytical and machine learning applications.
}
\newcommand{\conninetyfour}{
In this assignment, a Python program was developed to perform left, right, inner, and outer joins between two Pandas DataFrames. The program combines data from two tables based on a common key while applying different join strategies.

The implementation demonstrates the use of Pandas merge operations to retrieve matching and non-matching records according to the selected join type. A left join retains all records from the left DataFrame, a right join retains all records from the right DataFrame, an inner join returns only matching records, and an outer join includes all records from both DataFrames.

The exercise highlights the importance of data integration and relational data processing in data analysis. Understanding different join operations helps in combining datasets effectively and extracting meaningful information from multiple data sources.

This exercise improved knowledge of Pandas DataFrame manipulation and strengthened understanding of database-style join operations in Python.

Overall, the objective of performing left, right, inner, and outer joins between two DataFrames was successfully achieved using Pandas, demonstrating its effectiveness in managing and integrating structured data for analysis.
}
\newcommand{\conninetyfive}{
In this assignment, a Python program was developed to create a rolling window standard deviation column in a Pandas DataFrame. The program calculates the standard deviation of values within a specified moving window and stores the results in a new column.

The implementation demonstrates the use of Pandas rolling() and std() functions to perform window-based statistical analysis efficiently. Each value in the resulting column represents the variability of data within the defined rolling window.

The exercise highlights the importance of rolling statistics in time-series analysis, financial modeling, and trend monitoring. Rolling standard deviation is commonly used to measure fluctuations and assess the stability of data over time.

This exercise improved knowledge of Pandas DataFrame operations and strengthened understanding of statistical analysis and moving-window computations in Python.

Overall, the objective of creating a rolling window standard deviation column was successfully achieved using Pandas, demonstrating its effectiveness in performing advanced analytical and time-series processing tasks.
}
\newcommand{\conninetysix}{
In this assignment, a Python program was developed to create multiple layered charts in Altair with shared axes. The program combines two or more chart types, such as line charts, scatter plots, or area charts, into a single visualization while maintaining common x-axis and y-axis scales.

The implementation demonstrates the use of Altair's layering functionality, which allows multiple datasets or visual representations to be overlaid on the same coordinate system. Shared axes ensure consistent comparison and improve the readability of the combined visualization.

The exercise highlights the importance of advanced data visualization techniques in exploratory data analysis. Layered charts provide a comprehensive view of data by presenting multiple perspectives within a single graphical representation.

This exercise improved knowledge of Altair chart composition and strengthened understanding of layered visualizations, axis sharing, and data presentation techniques in Python.

Overall, the objective of creating multiple layered charts with shared axes was successfully achieved using Altair, demonstrating its flexibility and effectiveness for building rich and informative data visualizations.
}
\newcommand{\conninetyseven}{
In this assignment, a Python program was developed to calculate percentile values of a NumPy array. The program processes a set of numerical data and determines the values below which a specified percentage of observations fall.

The implementation demonstrates the use of NumPy's percentile() function to efficiently compute different percentile measures such as the 25th, 50th, and 75th percentiles. These statistical values provide useful insights into the distribution and spread of the dataset.

The exercise highlights the importance of percentile analysis in statistics, data science, and exploratory data analysis. Percentiles help summarize data, identify trends, and understand the relative standing of observations within a dataset.

This exercise improved knowledge of NumPy's statistical functions and strengthened understanding of descriptive statistics and data distribution concepts in Python.

Overall, the objective of calculating percentile values of a NumPy array was successfully achieved using NumPy, demonstrating its effectiveness in performing statistical and numerical computations.
}
\newcommand{\conninetyeight}{
In this assignment, a Python program was developed to create a time-series line chart with tooltips using the Altair visualization library. The chart plots data values against time, allowing trends and patterns to be observed over a specified period.

The implementation demonstrates the use of Altair's line chart functionality along with interactive tooltips. When the user hovers over data points, the tooltip displays detailed information such as the date, time, and corresponding values, enhancing the interactivity of the visualization.

The exercise highlights the importance of time-series analysis and interactive data visualization in understanding temporal trends and making data-driven decisions. Tooltips improve user experience by providing additional context without cluttering the chart.

This exercise improved knowledge of Altair visualization techniques and strengthened understanding of time-series data representation and interactive chart design in Python.

Overall, the objective of creating a time-series line chart with tooltips was successfully achieved using Altair, demonstrating its effectiveness in building informative and interactive visualizations for data analysis.
}
\newcommand{\conninetynine}{
In this assignment, a Python program was developed to split a Pandas DataFrame into training and testing datasets using random sampling. The program randomly divides the original dataset into two subsets, ensuring that the data is distributed fairly for model development and evaluation.

The implementation demonstrates the use of random sampling techniques to create separate training and testing sets. The training set is used to build and learn the model, while the testing set is used to evaluate its performance on unseen data.

The exercise highlights the importance of data partitioning in machine learning and predictive analytics. Proper separation of training and testing data helps prevent overfitting and provides a reliable assessment of model accuracy and generalization capability.
}
\newcommand{\cononehundred}{
In this assignment, a Python program was developed to load a CSV file using the Pandas library, clean missing data, compute summary statistics with NumPy, and visualize the results using Altair. The program performs a complete data analysis workflow, starting from data acquisition and ending with graphical representation.

The implementation demonstrates the use of Pandas for reading CSV files and handling missing values, NumPy for calculating statistical measures such as mean, median, standard deviation, and other summary statistics, and Altair for creating informative visualizations.

The exercise highlights the importance of data preprocessing, statistical analysis, and visualization in extracting meaningful insights from raw datasets. Cleaning missing values improves data quality, statistical summaries provide an overview of the dataset, and visualizations help reveal patterns and trends effectively.

This exercise improved knowledge of data analysis libraries in Python and strengthened understanding of the end-to-end data processing pipeline. It also enhanced skills in handling real-world datasets and presenting analytical results visually.

Overall, the objective of loading, cleaning, analyzing, and visualizing CSV data was successfully achieved using Pandas, NumPy, and Altair, demonstrating their combined effectiveness for data science and analytical tasks.
}